% =====================================================================================
\section{Discussion}
\label{sec:discussion}

% Paragraphs of the Discussion; shared by sec_discussion.tex and short_app_discussion.tex.
\paragraph{Without instructions: what the results show.} Every method without instructions loses at least the conflict of the inputs of two tasks (Proposition~\ref{prop:noinstructionprice}); that part is forced, and it is zero when no input occurs in two tasks. The rest is not forced. Without their instructions the two published methods that we ran changed \AudSeedsGainloraChanged{} and \AudSeedsInfloraChanged{} of \AudSeedsCompared{} answers. A recogniser on the frozen model's own states avoids most of this. It blocks no task and does not depend on the order of arrival. An audit estimates, from held-out inputs, the share of answers that change, and its variant \handicap{} bounds in advance what later tasks can take. On the 15 benchmark tasks, development data, the recogniser gives \AudDevLLamOneEnd{}\% of the test inputs to their own task; on 30 fresh tasks it gives \AudFreshFourthLEnd{}\% and passed the line fixed for it, as did two related rules on other fresh tasks; two tasks built on the same texts split their inputs under it (\AudScalePairOlderLEnd{}\% and \AudScalePairNewerLEnd{}\% reach their own task); \AudConfPairShare{}\% of their inputs occur word for word in both, and on those no method without instructions can do better than split them. On five tasks with trained answers (three seeds of one sequence, taken from the 15 development tasks), plain adapters behind the recogniser beat the two published methods without their definitions by a line fixed before the run, but there most of the gap is learning: they had a tenth of their published epochs. With their published number of epochs (one seed) they learn the five tasks as well as plain adapters do (\AudFullInfloraLearned{} and \AudFullGainloraLearned{} right after learning, against \AudOursRidgeLearned{}) and then forget \AudFullInfloraFT{} and \AudFullGainloraFT{} points, so there the gap is forgetting; forgetting grew with what had been learned. With the same 100 epochs plain adapters behind the recogniser end at \AudFullOursRidgeAP{} with \AudFullOursRidgeChanged{} of \AudFullCompared{} changed answers, the same routes and the same number of changed answers as with 10 epochs. The same run with their definitions gives \AudFullDefOriginalInfloraFT{} and \AudFullDefOriginalGainloraFT{} points of forgetting: in this setting the instruction is what held their answers, and our rules do not need it. It covers two methods on one sequence. The handicap is loose (\AudDevHPrimeTTwoBound{}\% against \AudDevHPrimeTTwoChanged{}\% measured). On fifteen tasks without definitions (one seed, our rules only) plain adapters behind the recogniser end \AudFifteenNoDefRidgeBelowOwn{} points below the ceiling with \AudFifteenNoDefRidgeChanged{} of \AudFifteenNoDefCompared{} changed answers. What remains: the published setting, the published methods without instructions on more than five tasks, a tight bound in advance, rewordings written by people, and tasks that share inputs.

\paragraph{Two rules, two kinds of statement.} The zones give an exact statement for the inputs inside them and none outside, and they read the instruction. The recogniser gives a measured statement for all inputs of a task (the audit), and its variant \handicap{} a loose one in advance; both read only the example. With instructions, on five tasks and three seeds, the two behaved as these statements say: the zones changed no answer and lost \AudUThreeHeavyCertBelowOwn{} points on average under the heavy rewording through misrouted inputs, and the recogniser changed at most \AudUThreeMaxChangedPerSeed{} of \AudUCompared{} answers per seed at every wording and stayed within \AudUThreeRidgeBelowOwnHigh{} points of the ceiling (Section~\ref{sec:exampleroutes}). Using the zone where an input is certified and the recogniser elsewhere gave the better of the two at the original, light and heavy wordings and not at the medium one (seed 42; Table~\ref{tab:instanceroutes}). Which rule a system should use depends on what it needs: a proof for unedited instructions, or routes that do not depend on the wording.

\paragraph{The third rule: zones from the inputs of a task.} Theorem~\ref{thm:bound} gives a guarantee that needs no instruction either: when each zone is calibrated on held-out inputs of its task, each input of a task changes its answer with probability at most a chosen level, for any model and any later tasks. The first tests (Section~\ref{sec:noinstruction}: routes on benchmark files without definitions, and a small trained model) met the bound and showed its price: with frozen word-vector features, most tasks are blocked at admission, because the zones of earlier tasks hold most of their inputs, and those tasks are then not learned. The bound moves the difficulty to one thing that is measured when a task arrives, the share of a task's inputs inside older zones. Without an instruction, telling tasks apart was what remained, and the recogniser (Sections~\ref{sec:audit} and~\ref{sec:recogniser}) takes it up.

\paragraph{What the comparison with the published methods shows, with instructions.} With the
instructions in the inputs, in our reduced setting the
certified switch ends far above InfLoRA and GainLoRA, and it would be easy to read this as ``they
forget and we do not''. That reading is wrong on five tasks. There both published methods have forgetting scores near zero (though \FiveInfChanged{} and \FiveGainChanged{} of \FiveInfCompared{} answers change). They had
learned less: right after learning they are at the level of guessing on the two tasks with two
labels. On fifteen tasks InfLoRA forgets \FifteenInfFT{} points and GainLoRA \FifteenGainFT{}, and
both had again learned less right after each task (\FifteenInfLearned{} and \FifteenGainLearned{}
against \FifteenCertLearned{}), so there the gap has a part from learning and a part from
forgetting. We did not find out why. Restricting an adapter to a subspace, or training only one of its
two matrices, may slow learning; ten epochs may simply be too few; and our set-up replaced one step
of their code by a substitute. Plain adapters learn these tasks well and then lose them; the switch
removes the loss without touching the learning. What our experiments support is therefore narrow:
switching protects without changing how an adapter is trained, and its price is paid where a new task collides with an old zone and, at scale, in routing: its fallback gave fewer machine-paraphrased inputs to their own task than a plain nearest-signature rule. They do not support a statement about how well the published
methods retain what they learn, because in the five-task runs there was little to retain, the
15-task run has one seed, and these methods were run with their definitions at their full
number of epochs for one seed only (Section~\ref{sec:nodeftrained}), and not with a larger model. Without their definitions the same two methods did hold their answers less well (Section~\ref{sec:nodefpublished}), with the same limits.

\paragraph{What the zones add to a nearest-task rule, and what they cost.} On the benchmark
sequences \nearest{} and \cert{} give the same routes and the same scores. This is expected: the
tasks of a benchmark are different from each other, and no later task tries to take an earlier
task's inputs. The two rules separate in two situations, and in opposite directions. When a later
task resembles an old one and inputs of the old task arrive in other words, the best-match rule
can hand them to the newcomer without notice, and the zones keep those that lie inside a zone (Section~\ref{sec:stress}, a constructed case). When
a later task resembles an old one and all instructions are unedited, the best-match rule serves
both tasks, and the zones refuse the newcomer (Section~\ref{sec:many}); in a collection of
\ManyTasks{} real tasks this happened to \ManyCollidingPct{}\% of the tasks. On the machine paraphrases of \WidenTasksEval{} real tasks this case was rare and the two rules gave the same routes (Section~\ref{sec:paraphrase}); there the zones added the proof for the inputs inside them, not better or more stable routing, and a plain nearest-signature rule gave more paraphrased inputs to their own task. We therefore cannot say that a system that expects reworded instructions is better served by the zones, only that it gets a proof for part of them. A system that expects
families of sibling tasks with fixed instructions is better served by the best match, or by zones
with a smaller radius, or by the variant of Section~\ref{sec:exactfirst}, which serves a sibling
and gives up the guarantee only for inputs that carry the sibling's instruction word for word. What
the zones offer in both situations, and a best-match rule does not, is that a conflict is known at
the moment the new task arrives, as the count of blocked training inputs, and that a certified
answer is final.

\paragraph{Exactness cuts both ways.} A certified answer never gets worse and never gets better.
InfLoRA and GainLoRA show small negative forgetting on some tasks: an old sentiment task improves
when a second sentiment task is learned. The switch cannot profit from this. One can still learn an
improved version of an old task as a new task, but under the same instruction its inputs carry the
old opening, so the old adapter keeps answering them under every zone rule of this paper; under
\recog{} the two versions would split such inputs (Proposition~\ref{prop:noinstructionprice}).
Replacing a
claim, with an explicit decision to give up the guarantee for one zone, is easy to add and is the
natural way to correct a task; we have not studied it.

\paragraph{Why averaged word-vectors are enough here, and when they are not.} The zone of a task is
built from its instruction. Two instructions that share few content words are far apart even after
frequent words are discounted, and an edit that keeps most content words stays close. That is why
light rewordings are certified. A heavy rewording can leave only a few content words, and then two
tasks about the same subject become hard to tell apart: in our runs the only inputs that the zones route wrongly
are tweets under a one-sentence instruction about positive and negative mood, which are given to
the other sentiment task. We tried features that read meaning rather than words: hidden states of a language model, a
frozen sentence encoder and a frozen inference model that judges whether two instructions ask for
the same thing (Section~\ref{sec:paraphrase}, Appendix~\ref{app:widen}). They raise the share of correct certificates, but none of them brings the wrong certificates under the 5\% we had asked for:
the lowest value, \WidenEncRulerTwentyFrM{}\% with \WidenEncRulerTwentyFrP{}\% correct certificates (\WidenRefTfiveFrM{}\% and \WidenRefTfiveFrP{}\% for \cert{}),
belongs to a setting singled out after the runs, and its advantage was smaller than the margin we fixed on held-out paraphrases. In the collisions we
inspected, sibling tasks that differ in one decisive word are put together as well. The theorem does not depend on the feature, only on the zones
being fixed when the task arrives and on membership being computed reproducibly by frozen parts (Remark~\ref{rem:otherzones}).

\paragraph{What the mechanical edits add to this picture.} The edit test (Section~\ref{sec:edits})
sharpens three points. First, robustness is a matter of length: an instruction of a hundred tokens
keeps its route under sixteen typos, an instruction of ten words loses it under eight. A system
built on this switch should give its tasks instructions of a few sentences. Second, the switch reads a bag of words. Shuffled
sentences are certified, and a sentence removed from the end matters only if it carried the words
that set the task apart. Third, typos are worse than missing words, because a misspelt word is
split into rare word pieces that the weighting of Eq.~\eqref{eq:weights} takes for content. A
spelling-tolerant feature might repair exactly this case, and by the argument above it could be
swapped in without touching the guarantee; we did not test one on typos.

\paragraph{What to do with an uncertified input.} The certificate separates two kinds of answers,
and a system can treat them differently. For a certified input the answer is final. For an
uncertified one the system knows that it is guessing the task, and it knows how good the guess is
likely to be, since $\min_k \delta_k(x)$ is the distance to the nearest zone in units of its radius (we did not test how well this distance predicts a correct route).
It can then answer anyway, as in our experiments, or ask for the instruction in its registered
form, or refuse. We did not study such policies; the edit test suggests that they matter for badly
damaged instructions, where the route is wrong for more than half of the inputs, all of which are
flagged.

\paragraph{The benchmark hands every method the task.} Every input of SuperNI begins with its task
definition, so the identity of the task can be read off the input by string matching. Methods that
are described as needing no task identity at test time therefore receive it as part of the text.
The exact key makes this explicit and reaches the full learning of plain adapters with no
forgetting, and it collapses to plain adapters when a single word of the instruction changes.
Evaluations of continual instruction tuning should include reworded instructions; otherwise a
lookup table of instructions is a perfect method. These benchmarks, with unedited instructions, cannot tell retention from recognition: a test of retention has to reword the instruction or take it away (Sections~\ref{sec:paraphrase} and~\ref{sec:nodefpublished}).

\paragraph{Evaluation noise.} The official loader scores an input against one reference answer
picked at random on every load. Identical answers then receive different scores in different
runs, about one point on a task with several references. On the 15 tasks it gave the certified switch a forgetting of \FifteenCertOffFT{} although none of
\FifteenCertCompared{} answers changed (Section~\ref{sec:fifteen}). When the claim is that nothing changed,
the answers should be compared, not the scores.


\section{Limitations}
\label{sec:limits}

\begin{itemize}
  \item \textbf{Reduced setting, and published methods that learned less.} T5-base instead of T5-large, and 10 epochs instead of 100 in all runs but a few, on free GPUs. In this setting the two published methods are at the level
  of guessing on two of the five tasks right after learning them, on the 15 tasks they learned less
  than plain adapters, and one step of their code was
  computed with a substitute for the library it uses, without a number-for-number check
  (Appendix~\ref{app:impl}). No conclusion about these methods in the published setting follows from our numbers, and the
  comparison without instructions (Section~\ref{sec:nodefpublished}) has the same limits: two methods, five tasks, a tenth of their epochs on three seeds and the full number on one.
  \item \textbf{Without instructions: two published methods, one sequence of five tasks, development
  tasks.} The comparison without instructions (Section~\ref{sec:nodefpublished}) covers two
  published methods and one sequence of five tasks with trained answers (three seeds), with a tenth
  of their published epochs. Plain adapters behind \recog{} or \handicap{} beat them by a line fixed
  before the run, but there most of the gap in the final score is learning, not retention. With their
  published number of epochs the two methods learned the tasks and then forgot \AudFullGainloraFTInt{}
  and \AudFullInfloraFTInt{} points (Section~\ref{sec:nodeftrained}); at the same number of epochs our rules end at \AudFullOursRidgeAPOne{} and \AudFullOursHandicapAPOne{}. With their definitions and the same number of epochs they forget \AudFullDefOriginalGainloraFTOne{} and \AudFullDefOriginalInfloraFTOne{} points, so in this setting the loss is due to the missing instruction. All of this is one seed. The rules were chosen on the 15 benchmark
  tasks, so their numbers there are development numbers (\AudDevLLamOneEnd{}\% of the test inputs at
  their own task after 15 tasks); the five GPU tasks are among them. The numbers on fresh tasks
  (\AudFreshFourthLEnd{}\% on 30 unseen tasks, \AudScaleLongThreeEnd{}\% on the second benchmark,
  \AudScaleAllHundredTwentyEnd{}\% with 120 tasks) are routes only. On fifteen tasks trained answers without instructions exist for our rules only, with one seed
  (Section~\ref{sec:nodeftrained}); there are none on the second benchmark, nor with the published methods given an oracle
  for the task or our recogniser in front.
  \item \textbf{The bound in advance is loose, and the split of fresh tasks is not exact.} The
  handicap gives \AudDevHPrimeTTwoBound{}\% on average against \AudDevHPrimeTTwoChanged{}\% measured; a
  tight bound in advance was not reached. The split of the fresh tasks does not ensure that test and
  held-out inputs are the same kind of draw (a later check with a random split moved the numbers by
  \AudScaleSplitMaxDiffEnd{} points or less).
  \item \textbf{Inputs that occur in two tasks, and rewordings by people.} Without instructions no
  method can do better than split inputs that occur in two tasks (two tasks built on the same texts
  reached \AudScalePairOlderLEnd{}\% and \AudScalePairNewerLEnd{}\% of their own inputs). Rewordings written by people were not
  tested, and the published setting (a larger model) was not reached.
  \item \textbf{Routes from the example alone: five tasks, and two fields.} The runs with
  instructions in which \recog{} sets the route from the example alone
  (Section~\ref{sec:exampleroutes}) have three seeds but only five tasks, and their rewordings were written by
  machine and are a development set. The benchmark gives the instruction and the example as two fields; a system that receives one
  string has to find where the example begins, and we did not test that.
  \item \textbf{Chosen tasks.} The five-task sequence was put together from tasks that one of the
  published methods had forgotten most in an earlier run of ours, and the unrelated inputs come
  from a benchmark from which we left out the sentiment tasks. The first choice was not varied; the
  second was varied for routes only (Section~\ref{sec:background}).
  \item \textbf{The closest method was run only as our own re-implementation.} Latent-LoRA, which
  also recognises tasks from frozen word-vectors, has no public code. We wrote it from its paper and
  ran it with a third of its epochs. It learned little, so we compare no scores. Its router gave
  every input to the right task and no answer changed. It was not run in the stress test
  (Appendix~\ref{app:latent}).
  \item \textbf{The scores do not separate the switch from simpler rules.} On the benchmark
  sequences a nearest-zone rule, and with unedited instructions even an exact match of the
  instruction, reach the same scores. What the certified switch adds there is the proof and the
  certificate, not points.
  \item \textbf{Few seeds, small test sets.} The five-task comparison has three seeds and the
  15-task comparison one. Test sets have about 100 inputs per task, so differences of a few points
  on one task are not meaningful, and two runs with the same seed already differ by a few tenths
  of a point.
  \item \textbf{At scale the switch adds the proof, not better routing.} On machine paraphrases of \WidenTasksEval{} real tasks the nearest-zone rule changed \RerouteNearNChanged{} of \RerouteInputs{} routes after the tasks arrived
  and the certified switch the same \RerouteCertNChanged{}, against the \RerouteExpectNearLow{}\% to \RerouteExpectNearHigh{}\% we had expected. With the fallback used in all GPU runs the switch routed worse at scale than a plain nearest-signature rule
  (\RerouteCertOwnEnd{}\% against \RerouteSigOwnEnd{}\% of the inputs to their own task). A variant of the fallback, scaled by each task's distance to unrelated text and tested at the end on routes only, is right on the benchmark (\RerouteBgHeavyFive{}\% under heavy rewording) and at scale
  (\RerouteBgItOwn{}\% without zones and \RerouteCertFourItOwn{}\% with the zones in front, on paraphrases through a fourth language); even then the best rule without zones gives about five points more right routes, so the zones add the proof and not better routing. The two fallbacks trade places between the tests, and the variant has not been run with trained answers (Section~\ref{sec:paraphrase}).
  \item \textbf{The zones from the inputs of a task block later tasks.}
  The bound of Theorem~\ref{thm:bound} holds for test inputs that are like the training inputs of their task and was met in every run (\BoundAllChangedLow{}\% to \BoundAllChangedHigh{}\% of the routes changed against the \BoundAlpha{}\% allowed), but with frozen word-vector features most later tasks were blocked when they arrived (Section~\ref{sec:noinstruction}).
  \item \textbf{The zones are narrow at scale.} On machine paraphrases of \WidenTasksEval{} real tasks \cert{} certifies \WidenRefTfivePLow{}\% to \WidenRefTfivePHigh{}\% of the inputs for their own task and \WidenRefTfiveMLow{}\% to \WidenRefTfiveMHigh{}\% for another task,
  depending on the pivot language (Section~\ref{sec:paraphrase}); about three quarters reach their own task, against \FiveCertRightLight{}\%, \FiveCertRightMedium{}\% and \FiveCertRightHeavy{}\% on the five benchmark tasks, which are far apart. None of five attempts
  to widen the zones (hidden states, a paraphrase ruler, a sentence encoder with a conformal radius, an inference model as judge, the same judge sentence by sentence) met the bar set for it before the runs: the share of correct certificates rose, but in every attempt the wrong certificates rose with it
  (hidden states, paraphrase ruler) or stayed above the bar set for them (sentence encoder \WidenEncPrimM{}\%, judge \WidenJudgeFrM{}\%, sentence-by-sentence judge \WidenSentItM{}\% with \WidenSentItP{}\% correct certificates on the Italian pivot, against 9\% for a failure and 5\% for a success)
  (Appendix~\ref{app:widen}). Machine paraphrases are not rewordings by people. A sentence-encoder zone chosen after the runs was ahead of \cert{}
  on each measure on held-out paraphrases of the same tasks, but by less than the margin fixed in advance, and not on every measure in other random splits, so we do not claim it as an improvement.
  \item \textbf{The reworded instructions are a development set.} They were written by machine during
  the development of the method, they share words with the originals, and two design
  choices (the radius factor, and the rule for inputs outside the zones together with the word
  weights) were made while looking at where they were routed. An independent test needs rewordings
  written by people who have not seen the method.
  \item \textbf{The zone guarantee covers certified inputs only.} With medium and heavy rewording most
  inputs are outside the zones. They kept their answers in our runs, but nothing guarantees it, and
  the stress test shows that a later look-alike task can take them. How many of them reach their
  own task also depends on the unrelated inputs that set the scale of distances: between
  \BackgroundHeavyLow{}\% and \BackgroundHeavyHigh{}\% under heavy rewording
  (Section~\ref{sec:background}).
  \item \textbf{Uncertified routes fail under heavy damage.} In the mechanical edit test, sixteen
  typos in the instruction leave only 45\% of the inputs on their own task, and tasks with very
  short instructions lose their route after a few edits. The method flags these inputs; it does not
  rescue them.
  \item \textbf{The switch ignores word order.} Its feature is a weighted average of the first
  tokens, so two instructions made of the same words in a different order cannot be told apart.
  \item \textbf{Several tests are about routes only.} The machine-paraphrase test, the re-routing study, the five attempts to widen the zones, the Spanish check, the stress test, the mechanical edit test, the test of task orders and the test with other embedding tables compute where inputs are sent; they do not measure what a model answers once an input has reached its adapters. The paraphrase tests use one split and one arrival order (except the five extra splits of the Spanish check).
  \item \textbf{The instruction zones need an instruction.} The zones of Section~\ref{sec:zones} are
  built from the common opening of a task's inputs; a task whose inputs share no opening has no
  signature. Such tasks are served by \recog{} (Sections~\ref{sec:recogniser}
  and~\ref{sec:nodeftrained}), not by the zones.
  \item \textbf{Older tasks have priority, and in a large collection this is costly.} A new task
  that collides with an old zone is blocked on the colliding inputs. Among \ManyTasks{} real tasks
  this happened to \ManyCollidingPct{}\% of the tasks, almost all of them siblings of an older task
  with a nearly identical instruction, and \ManyAllHeld{} tasks could not be learned at all under
  their given instruction. The collision is reported, but the zones offer no remedy except
  rewording the new task's instruction, a smaller radius factor, or giving up the old claim; the
  variant \exactfirst{} serves such tasks with a weaker guarantee and was tested on routes only. On
  unedited instructions a nearest-zone rule without frozen zones gives more inputs to their own
  task (\ManyNearRightPct{}\% against \ManyCertRightPct{}\%).
  \item \textbf{No backward transfer.} Certified answers cannot improve through later tasks.
  \item \textbf{One trained model family.} Training and answers were tested with T5-base and, at
  small scale on a CPU, with the stand-alone code on small models. For the other models only the
  routes were checked.
  \item \textbf{Memory grows with the number of tasks.} One adapter per task is kept, as in all
  methods of this family.
\end{itemize}

\section{Conclusion}
\label{sec:conclusion}

If the adapters of old tasks are frozen and each task keeps one set, an answer changes only if its
route changes, so forgetting is a question about the rule that sets the route. Without instructions
every method loses at least the conflict of the inputs of two tasks, and the published benchmarks
hide this, because their inputs begin with the instruction of the task. When it was removed, the
two published methods that we ran changed \AudSeedsGainloraChanged{} and \AudSeedsInfloraChanged{}
of \AudSeedsCompared{} answers, and this loss was not forced. Our main rule is a recogniser on the
frozen model's own states: it blocks no task, is the same for every order of arrival, and comes with
an audit of the share of changed answers and a bound in advance on what later tasks can take, which
is loose. It gave \AudDevLLamOneEnd{}\% of the test inputs of the 15 development tasks to their own
task, \AudScaleLongThreeEnd{}\% on a second benchmark, \AudFreshFourthLEnd{}\% on 30 unseen tasks and
\AudScaleAllHundredTwentyEnd{}\% with 120 tasks. In a reduced comparison on five tasks, plain
adapters behind it beat the two published methods without their instructions by a line fixed before
the run. With a tenth of their epochs most of that gap is learning; with their published number of
epochs (one seed) they learned the tasks and then forgot \AudFullGainloraFTInt{} and
\AudFullInfloraFTInt{} points, ending at \AudFullGainloraAPOne{} and \AudFullInfloraAPOne{} against
\AudFullOursRidgeAPOne{} for plain adapters behind the recogniser with the same number of epochs. The five tasks are among the development tasks.

The recogniser reads only the example, so its routes are the same for every wording of the
instruction: with instructions (five tasks, three seeds, machine-written rewordings)
plain adapters behind it ended at \AudUThreeHeavyRidgeAPOne{} under a heavy rewording
where the zones ended at \AudUThreeHeavyCertAPOne{}, and changed at most \AudUThreeMaxChangedPerSeed{} of
\AudUCompared{} answers per seed where the zones changed none.

When inputs do carry an instruction, an exact statement is possible: each task claims a zone built
from the frozen word-vectors of its instruction, and the oldest claim wins. The answer to an input
inside a zone then cannot change, for any model and any later tasks, and the price, a new task that
collides with an old claim, is visible when it arises. In a reduced setting on a real model, on five
tasks and on 15, this rule kept every answer of the old tasks while learning as much as unprotected
adapters, with the original instructions and, in the order we ran, with reworded ones. Simpler rules
that recognise the task reach the same scores there; what the zones add is the proof. In a
constructed case, a look-alike task re-routed certified old inputs under a best-match rule but not
under the zones (routes only). On machine paraphrases of real tasks, however, a nearest-zone rule
changed the same \RerouteNearNChanged{} of \RerouteInputs{} routes. A plain nearest-signature rule
gave more of them to their own task (\RerouteSigOwnEnd{}\% against \RerouteCertOwnEnd{}\%) but did
much worse on the benchmark rewordings. A variant of the fallback, fixed at the end, does well in
both places on route tests, yet the best rule without zones still routes about five points better. The limits of the zones are equally
plain: the guarantee ends at the border of the zones, the zones are built from words and not from
meaning (on machine paraphrases of real tasks they certify \WidenRefTfivePLow{}\% to \WidenRefTfivePHigh{}\% of the inputs for their own task, and none of five attempts to widen
them met the bar set for it in advance), a mechanical test showed that the route of an input outside the zones is lost when its
instruction is badly damaged, and in a large collection of real tasks about one task in six arrived
inside an older zone and was blocked. The version of the zones built from the inputs of a task has a proved bound on forgetting, met in every run, but with frozen word-vector features it blocks most later tasks, which are then not learned.

What is missing: the larger model of the published setting, more than one seed at the full number of epochs, the published methods without instructions on more than five tasks, a
tight bound in advance, rewordings written by people, and tasks that share inputs, which no method without
instructions can do better than split. All of this was measured with a small model, and all but a few runs with a short training budget.

